%%%%%%%%%%%%%%%%
%%% deut 30 %%%
%%%%%%%%%%%%%%%%

\Chapter{30}

\Verse{1}
{
	\hDtrB{וְהָיָה כִי יָבֹאוּ עָלֶיךָ כּל הַדְּבָרִים הָאֵלֶּה הַבְּרָכָה וְהַקְּלָלָה אֲשֶׁר נָתַתִּי לְפָנֶיךָ וַהֲשֵׁבֹתָ אֶל לְבָבֶךָ בְּכל הַגּוֹיִם אֲשֶׁר הִדִּיחֲךָ יְהֹוָה אֱלֹהֶיךָ שָׁמָּה׃}
}
{
	\eDtrB{
		30 And it will be, when all these things, the blessing and
		the curse that I’ve put in front of you, will come upon
		you, and you’ll store it in your heart among the nations
		to which YHWH, your God, has driven you,
	}
}
{
%	\footnote{
%		your God has driven you. The word “driven” is ironic here. It has been
%		used three times until now, always referring to persons who would drive
%		Israelites away from their God (Deut 13:6,11,14). But now it describes
%		how YHWH drives Israel away to other lands because of Israel’s apostasy.
%		The term is thus another allusion to the fact of divine punishment to
%		fit the crime, which is to say: there is justice.
%	}
}

\Verse{2}
{
	\hDtrB{וְשַׁבְתָּ עַד יְהֹוָה אֱלֹהֶיךָ וְשָׁמַעְתָּ בְקֹלוֹ כְּכֹל אֲשֶׁר אָנֹכִי מְצַוְּךָ הַיּוֹם אַתָּה וּבָנֶיךָ בְּכל לְבָבְךָ וּבְכל נַפְשֶׁךָ׃}
}
{
	\eDtrB{
		and you’ll come back to YHWH, your God, and listen to His
		voice, according to everything that I command you today,
		you and your children, with all your heart and with all
		your soul,
	}
}
{
}

\Verse{3}
{
	\hDtrB{וְשָׁב יְהֹוָה אֱלֹהֶיךָ אֶת שְׁבוּתְךָ וְרִחֲמֶךָ וְשָׁב וְקִבֶּצְךָ מִכּל הָעַמִּים אֲשֶׁר הֱפִיצְךָ יְהֹוָה אֱלֹהֶיךָ שָׁמָּה׃}
}
{
	\eDtrB{
		that YHWH, your God, will bring back your captivity and be
		merciful to you. And He’ll come back and gather you from
		all the peoples to which YHWH, your God, has scattered
		you.
	}
}
{
%	\footnote{
%		bring back your captivity. This expression occurs about twenty-five
%		times in the Tanak. It is grammatically unclear. The verb
%		([image][image]b) would normally be in the Hiphil, but here (and
%		usually) it is in the intransitive Qal. The noun may reflect the root
%		[image]wb, “bring back” (making it a cognate accusative with the verb),
%		or it may reflect the root [image]bh, “captive.” On the former root, the
%		phrase would mean “to produce a coming back.” On the latter root, it
%		would mean “to bring back a captivity.” These two meanings are very
%		close to each other in any case. Those who take this phrase to mean
%		figuratively “to restore your fortunes” have a good parallel in Job
%		42:10; but all the occurrences in Jeremiah, which is linked closely with
%		Deuteronomy, fit the meaning of “coming back” or “captivity” far better.
%		Note especially Jer 30:16–23 and 48:46–47, in which this expression
%		occurs in proximity to the related word for captivity: [image]ebî.
%	}
}

\Verse{4}
{
	\hDtrB{אִם יִהְיֶה נִדַּחֲךָ בִּקְצֵה הַשָּׁמָיִם מִשָּׁם יְקַבֶּצְךָ יְהֹוָה אֱלֹהֶיךָ וּמִשָּׁם יִקָּחֶךָ׃}
}
{
	\eDtrB{
		If you’ll be driven to the end of the skies, YHWH, your
		God, will gather you from there, and He’ll take you from
		there.
	}
}
{
}

\Verse{5}
{
	\hDtrB{וֶהֱבִיאֲךָ יְהֹוָה אֱלֹהֶיךָ אֶל הָאָרֶץ אֲשֶׁר יָרְשׁוּ אֲבֹתֶיךָ וִירִשְׁתָּהּ וְהֵיטִבְךָ וְהִרְבְּךָ מֵאֲבֹתֶיךָ׃}
}
{
	\eDtrB{
		And YHWH, your God, will bring you to the land that your
		fathers possessed, and you’ll take possession of it, and
		He’ll be good to you and multiply you more than your
		fathers.
	}
}
{
}

\Verse{6}
{
	\hDtrB{וּמָל יְהֹוָה אֱלֹהֶיךָ אֶת לְבָבְךָ וְאֶת לְבַב זַרְעֶךָ לְאַהֲבָה אֶת יְהֹוָה אֱלֹהֶיךָ בְּכל לְבָבְךָ וּבְכל נַפְשְׁךָ לְמַעַן חַיֶּיךָ׃}
}
{
	\eDtrB{
		And YHWH, your God, will circumcise your heart and your
		seed’s heart so as to love YHWH, your God, with all your
		heart and with all your soul so that you’ll live.
	}
}
{
%	\footnote{
%		God will circumcise your heart. Earlier, Moses directs the people to
%		circumcise their hearts themselves (see Deut 10:16 and the comment
%		there), but now he speaks of God circumcising their hearts for them. The
%		difference is that now he is speaking of a future condition: the people
%		will have violated the covenant, suffered the curses, and turned back to
%		their God; and God, in response, will bring them back and then will
%		circumcise their hearts. This in turn will enable them to love the deity
%		with all their hearts and souls—which Moses had commanded them to do
%		earlier as well. The divine-human relationship is pictured as mutual.
%		Humans require the experience of being away from their God, on their
%		own. And then, more experienced, more understanding, having suffered and
%		grown, humans look to God anew. And then the new encounter with the
%		divine transforms them. Footnote 30:6. your seed’s heart. Meaning: the
%		hearts of your progeny.
%	}
}

\Verse{7}
{
	\hDtrB{וְנָתַן יְהֹוָה אֱלֹהֶיךָ אֵת כּל הָאָלוֹת הָאֵלֶּה עַל אֹיְבֶיךָ וְעַל שֹׂנְאֶיךָ אֲשֶׁר רְדָפוּךָ׃}
}
{
	\eDtrB{
		And YHWH, your God, will put all these curses on your
		enemies and on those who hate you who have pursued you.
	}
}
{
}

\Verse{8}
{
	\hDtrB{וְאַתָּה תָשׁוּב וְשָׁמַעְתָּ בְּקוֹל יְהֹוָה וְעָשִׂיתָ אֶת כּל מִצְוֺתָיו אֲשֶׁר אָנֹכִי מְצַוְּךָ הַיּוֹם׃}
}
{
	\eDtrB{
		And you’ll come back and listen to YHWH’s voice and do all
		His commandments that I command you today.
	}
}
{
}

\Verse{9}
{
	\hDtrB{וְהוֹתִירְךָ יְהֹוָה אֱלֹהֶיךָ בְּכֹל מַעֲשֵׂה יָדֶךָ בִּפְרִי בִטְנְךָ וּבִפְרִי בְהֶמְתְּךָ וּבִפְרִי אַדְמָתְךָ לְטֹבָה כִּי יָשׁוּב יְהֹוָה לָשׂוּשׂ עָלֶיךָ לְטוֹב כַּאֲשֶׁר שָׂשׂ עַל אֲבֹתֶיךָ׃}
}
{
	\eDtrB{
		And YHWH, your God, will give you extra of all your hand’s
		work, of the fruit of your womb and of the fruit of your
		animals and the fruit of your land for good. Because YHWH
		will come back to have satisfaction over you for good as
		He had satisfaction over your fathers,
	}
}
{
}

\Verse{10}
{
	\hDtrB{כִּי תִשְׁמַע בְּקוֹל יְהֹוָה אֱלֹהֶיךָ לִשְׁמֹר מִצְוֺתָיו וְחֻקֹּתָיו הַכְּתוּבָה בְּסֵפֶר הַתּוֹרָה הַזֶּה כִּי תָשׁוּב אֶל יְהֹוָה אֱלֹהֶיךָ בְּכל לְבָבְךָ וּבְכל נַפְשֶׁךָ׃}
}
{
	\eDtrB{
		when you’ll listen to the voice of YHWH, your God, to
		observe His commandments and His laws, written in this
		scroll of instruction, when you’ll come back to YHWH, your
		God, with all your heart and with all your soul.
	}
}
{
}

\Verse{11}
{
	\hDtrA{כִּי הַמִּצְוָה הַזֹּאת אֲשֶׁר אָנֹכִי מְצַוְּךָ הַיּוֹם לֹא נִפְלֵאת הִוא מִמְּךָ וְלֹא רְחֹקָה הִוא׃}
}
{
	\eDtrA{
		“Because this commandment that I command you today: it’s
		not too wondrous for you, and it’s not too far.
	}
}
{
%	\footnote{
%		it’s not too wondrous for you. Now Moses finishes what he started to say
%		ten verses earlier about the hidden things belonging to God and the
%		revealed things belonging to humans (29:28). In one of the most
%		beautiful passages in the Torah (vv. 11–14), he explains what this
%		meant: The commandments are not enigmatic, they do not reside in a
%		distant realm, and they do not require an intermediary. They are already
%		made known. And they are within a human’s ability to do.
%	}
}

\Verse{12}
{
	\hDtrA{לֹא בַשָּׁמַיִם הִוא לֵאמֹר מִי יַעֲלֶה לָּנוּ הַשָּׁמַיְמָה וְיִקָּחֶהָ לָּנוּ וְיַשְׁמִעֵנוּ אֹתָהּ וְנַעֲשֶׂנָּה׃}
}
{
	\eDtrA{
		It’s not in the skies, that one would say, ‘Who will go up
		for us to the skies and get it for us and enable us to
		hear it so we’ll do it?’
	}
}
{
}

\Verse{13}
{
	\hDtrA{וְלֹא מֵעֵבֶר לַיָּם הִוא לֵאמֹר מִי יַעֲבר לָנוּ אֶל עֵבֶר הַיָּם וְיִקָּחֶהָ לָּנוּ וְיַשְׁמִעֵנוּ אֹתָהּ וְנַעֲשֶׂנָּה׃}
}
{
	\eDtrA{
		And it’s not across the sea, that one would say, ‘Who will
		cross for us, across the sea, and get it for us and enable
		us to hear it so we’ll do it?’
	}
}
{
}

\Verse{14}
{
	\hDtrA{כִּי קָרוֹב אֵלֶיךָ הַדָּבָר מְאֹד בְּפִיךָ וּבִלְבָבְךָ לַעֲשֹׂתוֹ׃}
}
{
	\eDtrA{
		But the thing is very close to you, in your mouth, and in
		your heart, to do it.
	}
}
{
}

\Verse{15}
{
	\hDtrB{רְאֵה נָתַתִּי לְפָנֶיךָ הַיּוֹם אֶת הַחַיִּים וְאֶת הַטּוֹב וְאֶת הַמָּוֶת וְאֶת הָרָע׃}
}
{
	\eDtrB{
		“See: I’ve put in front of you today life and good, and
		death and bad,
	}
}
{
}

\Verse{16}
{
	\hDtrB{אֲשֶׁר אָנֹכִי מְצַוְּךָ הַיּוֹם לְאַהֲבָה אֶת יְהֹוָה אֱלֹהֶיךָ לָלֶכֶת בִּדְרָכָיו וְלִשְׁמֹר מִצְוֺתָיו וְחֻקֹּתָיו וּמִשְׁפָּטָיו וְחָיִיתָ וְרָבִיתָ וּבֵרַכְךָ יְהֹוָה אֱלֹהֶיךָ בָּאָרֶץ אֲשֶׁר אַתָּה בָא שָׁמָּה לְרִשְׁתָּהּ׃}
}
{
	\eDtrB{
		in that I command you today to love YHWH, your God, to go
		in His ways and to observe His commandments and His laws
		and His judgments so you’ll live and multiply and bless
		YHWH, your God, in the land to which you’re coming to take
		possession of it.
	}
}
{
%	\footnote{
%		so you’ll live and multiply. This is a change of wording from the
%		beginning of the Torah, where the command is to “be fruitful and
%		multiply” (Gen 1:28). The command to grow in numbers is for the purpose
%		of life. As the human population grows to the point that its number is a
%		threat to life, humankind and its leaders need enormous wisdom to
%		determine what to do. The wording there in Genesis is: “And God blessed
%		them, and God said to them, ‘Be fruitful and multiply … ’” (Gen 1:28).
%		That first commandment of the Torah is thus presented not as a strict
%		order, but as a blessing. And here at the end of the Torah as well, the
%		text is: “so you’ll live and multiply and bless YHWH.” In the Torah, the
%		creator blesses humans with the opportunity to grow, and humans will
%		bless their creator for it once it is fulfilled. The Bible’s first
%		commandment is now becoming humankind’s greatest challenge.
%	}
}

\Verse{17}
{
	\hDtrB{וְאִם יִפְנֶה לְבָבְךָ וְלֹא תִשְׁמָע וְנִדַּחְתָּ וְהִשְׁתַּחֲוִיתָ לֵאלֹהִים אֲחֵרִים וַעֲבַדְתָּם׃}
}
{
	\eDtrB{
		But if your heart will turn away, and you won’t listen,
		and you’ll be driven so that you’ll bow to other gods and
		serve them,
	}
}
{
}

\Verse{18}
{
	\hDtrB{הִגַּדְתִּי לָכֶם הַיּוֹם כִּי אָבֹד תֹּאבֵדוּן לֹא תַאֲרִיכֻן יָמִים עַל הָאֲדָמָה אֲשֶׁר אַתָּה עֹבֵר אֶת הַיַּרְדֵּן לָבוֹא שָׁמָּה לְרִשְׁתָּהּ׃}
}
{
	\eDtrB{
		I’ve told you today that you’ll perish. You won’t extend
		days on the land to which you’re crossing the Jordan to
		come to take possession of it.
	}
}
{
}

\Verse{19}
{
	\hDtrB{הַעִדֹתִי בָכֶם הַיּוֹם אֶת הַשָּׁמַיִם וְאֶת הָאָרֶץ הַחַיִּים וְהַמָּוֶת נָתַתִּי לְפָנֶיךָ הַבְּרָכָה וְהַקְּלָלָה וּבָחַרְתָּ בַּחַיִּים לְמַעַן תִּחְיֶה אַתָּה וְזַרְעֶךָ׃}
}
{
	\eDtrB{
		I call the skies and the earth to witness regarding you
		today: I’ve put life and death in front of you, blessing
		and curse. And you shall choose life, so you’ll live, you
		and your seed,
	}
}
{
%	\footnote{
%		choose life. The theme of the path to life began early in Moses’ speech
%		(Deut 4:1) and culminates here in the last words of the speech. This
%		focus at the conclusion of the Torah returns us to the Torah’s opening:
%		the loss of the tree of life. Humans lose access to the tree of life as
%		the price of having gained access to the tree of knowledge of good and
%		bad. Now the people are told, “I’ve put in front of you life and good,
%		blessing and curse.” Using the knowledge of good and bad, and choosing
%		to do good, is the path back to life—not necessarily eternal life
%		(though who knows?), but meaningful life, fulfilling life. As the book
%		of Proverbs says about knowledge and wisdom: “It is a tree of life.” And
%		Jews sing this verse from Proverbs each Sabbath after reading the Torah
%		and returning it to the ark.
%	}
}

\Verse{20}
{
	\hDtrB{לְאַהֲבָה אֶת יְהֹוָה אֱלֹהֶיךָ לִשְׁמֹעַ בְּקֹלוֹ וּלְדבְקָה בוֹ כִּי הוּא חַיֶּיךָ וְאֹרֶךְ יָמֶיךָ לָשֶׁבֶת עַל הָאֲדָמָה אֲשֶׁר נִשְׁבַּע יְהֹוָה לַאֲבֹתֶיךָ לְאַבְרָהָם לְיִצְחָק וּלְיַעֲקֹב לָתֵת לָהֶם׃}
}
{
	\eDtrB{
		to love YHWH, your God, to listen to His voice and to
		cling to Him, because He is your life and the extension of
		your days to reside on the land that YHWH swore to your
		fathers, to Abraham, to Isaac, and to Jacob, to give to
		them.”
	}
}
{
%	\footnote{
%		to cling to Him. “Cling” is the same term that is used for man’s
%		attachment to woman at the beginning of the Torah (Gen 2:24). This is
%		part of a string of terms in this passage (vv. 15–20) that call to mind
%		the opening chapters of the Torah: life, death, good and bad, the skies
%		and earth, the ground, the day (cf. “in the day” and “in the day you eat
%		from it,” Gen 2:4,17), “your seed,” be fruitful, to give, to tell, to
%		command (cf. Gen 3:11), to watch (or observe, Hebrew [image]mr), and to
%		listen to the voice (of God; cf. “You listened to your woman’s voice,”
%		Gen 3:17). And these are preceded by references to fruit (Deut 30:9),
%		and to the skies and the sea (30:12–13). This cluster of language here
%		in the last sentences of Moses’ address forms a great connection back to
%		the start of Genesis and a reminder of the unity of the Torah. Footnote
%		30:20. He is your life. This can also be understood to mean “It is your
%		life.” Translators are split on this. The former means that Moses is
%		saying that God is the people’s life and the source of their being in
%		their land a long time. The latter means that the people’s choice to
%		love, listen, and cling to God is their life and the source of lengthy
%		time in the land. That is, the question is whether it is God or it is
%		people’s feeling about God that brings this about. I do not know which
%		is correct.
%	}
}
